\let\negmedspace\undefined
\let\negthickspace\undefined
\documentclass[journal]{IEEEtran}
\usepackage[a5paper, margin=10mm, onecolumn]{geometry}
%\usepackage{lmodern} % Ensure lmodern is loaded for pdflatex
\usepackage{tfrupee} % Include tfrupee package

\setlength{\headheight}{1cm} % Set the height of the header box
\setlength{\headsep}{0mm}     % Set the distance between the header box and the top of the text

\usepackage{gvv-book}
\usepackage{gvv}
\usepackage{cite}
\usepackage{amsmath,amssymb,amsfonts,amsthm}
\usepackage{algorithmic}
\usepackage{graphicx}
\usepackage{textcomp}
\usepackage{xcolor}
%\usepackage{txfonts}
\usepackage{listings}
\usepackage{enumitem}
\usepackage{mathtools}
\usepackage{gensymb}
\usepackage{comment}
\usepackage[breaklinks=true]{hyperref}
\usepackage{tkz-euclide} 
\usepackage{listings}
% \usepackage{gvv}                                        
\def\inputGnumericTable{}                                 
\usepackage[latin1]{inputenc}                                
\usepackage{color}                                            
\usepackage{array}                                            
\usepackage{longtable}                                       
\usepackage{calc}                                             
\usepackage{multirow}                                         
\usepackage{hhline}                                           
\usepackage{ifthen}                                           
\usepackage{lscape}
\usepackage{circuitikz}
\tikzstyle{block} = [rectangle, draw, fill=blue!20, 
    text width=4em, text centered, rounded corners, minimum height=3em]
\tikzstyle{sum} = [draw, fill=blue!10, circle, minimum size=1cm, node distance=1.5cm]
\tikzstyle{input} = [coordinate]
\tikzstyle{output} = [coordinate]


\begin{document}

\bibliographystyle{IEEEtran}
\vspace{3cm}

\title{5.7.3}
\author{AI25BTECH11030 -Sarvesh Tamgade}
{\let\newpage\relax\maketitle}

\renewcommand{\thefigure}{\theenumi}
\renewcommand{\thetable}{\theenumi}
\setlength{\intextsep}{10pt} 


\numberwithin{equation}{enumi}
\numberwithin{figure}{enumi}
\renewcommand{\thetable}{\theenumi}


\textbf{Question}: For the matrix
\[
\myvec{
1 & 1 & 1 \\
1 & 2 & -2 \\
2 & -1 & 3
},
\]
show that
\begin{align}
A^3 - 6A^2 + 5A + 11I &= 0.
\end{align}
Find \(A^{-1}\).


\textbf{Solution:} \
The characteristic equation is
\begin{align}
\left|\myvec{
1-\lambda & 1 & 1 \\
1 & 2-\lambda & -2 \\
2 & -1 & 3-\lambda
} \right| &= 0,
\end{align}
which expands to
\begin{align}
\lambda^3 - 6 \lambda^2 + 5 \lambda + 11 &= 0.
\end{align}
By Cayley-Hamilton theorem,
\begin{align}
A^3 - 6A^2 + 5A + 11I &= 0.
\end{align}
Rearranging,
\begin{align}
11I &= -A^3 + 6A^2 - 5A,
\end{align}
and multiplying both sides by \(A^{-1}\),
\begin{align}
A^{-1} &= \frac{1}{11}(-A^2 + 6A - 5I) \\
&= \frac{1}{11} \myvec{
-3 & 4 & 4 \\
7 & 0 & -3 \\
5 & -3 & 0
}.
\end{align}

\textbf{Final Answer:}
\begin{align}
A^{-1} &= \myvec{
-\frac{3}{11} & \frac{4}{11} & \frac{4}{11} \\
\frac{7}{11} & 0 & -\frac{3}{11} \\
\frac{5}{11} & -\frac{3}{11} & 0
}.
\end{align}
\end{document}  